\documentclass[a4paper, anonymous, UKenglish,cleveref, autoref, thm-restate]{lipics-v2021}

% Required packages from the original paper
\usepackage{mathpartir}
\usepackage{listings}
\usepackage{syntax}
\usepackage{algorithm}
\usepackage{algpseudocode}
\usepackage{mathtools}
\usepackage{amssymb}

% Define colors for code listings
\definecolor{dkgreen}{rgb}{0,0.6,0}
\definecolor{ltblue}{rgb}{0,0.4,0.4}
\definecolor{dkviolet}{rgb}{0.3,0,0.5}
\definecolor{dkblue}{rgb}{0,0,0.6}
\definecolor{dkred}{rgb}{0.6,0,0}
\definecolor{codegray}{rgb}{0.5,0.5,0.5}
\definecolor{codepurple}{rgb}{0.58,0,0.82}
\definecolor{backcolour}{rgb}{0.95,0.95,0.92}

% Coq language definition for listings
\lstdefinelanguage{Coq}{ 
    % Anything betweeen $ becomes LaTeX math mode
    mathescape=true,
    % Comments may or not include Latex commands
    texcl=false, 
    % Vernacular commands
    morekeywords=[1]{Section, Module, End, Require, Import, Export,
        Variable, Variables, Parameter, Parameters, Axiom, Hypothesis,
        Hypotheses, Notation, Local, Tactic, Reserved, Scope, Open, Close,
        Bind, Delimit, Definition, Let, Ltac, Fixpoint, CoFixpoint, Add,
        Morphism, Relation, Implicit, Arguments, Unset, Contextual,
        Strict, Prenex, Implicits, Inductive, CoInductive, Record,
        Structure, Canonical, Coercion, Context, Class, Global, Instance,
        Program, Infix, Theorem, Lemma, Corollary, Proposition, Fact,
        Remark, Example, Proof, Goal, Save, Qed, Defined, Hint, Resolve,
        Rewrite, View, Search, Show, Print, Printing, All, Eval, Check,
        Projections, inside, outside, Def},
    % Gallina
    morekeywords=[2]{forall, exists, exists2, fun, fix, cofix, struct,
        match, with, end, as, in, return, let, if, is, then, else, for, of,
        nosimpl, when},
    % Sorts
    morekeywords=[3]{Type, Prop, Set, true, false, option},
    % Various tactics
    morekeywords=[4]{pose, set, move, case, elim, apply, clear, hnf,
        intro, intros, generalize, rename, pattern, after, destruct,
        induction, using, refine, inversion, injection, rewrite, congr,
        unlock, compute, ring, field, fourier, replace, fold, unfold,
        change, cutrewrite, simpl, have, suff, wlog, suffices, without,
        loss, nat_norm, assert, cut, trivial, revert, bool_congr, nat_congr,
        symmetry, transitivity, auto, split, left, right, autorewrite},
    % Terminators
    morekeywords=[5]{by, done, exact, reflexivity, tauto, romega, omega,
        assumption, solve, contradiction, discriminate},
    % Control
    morekeywords=[6]{do, last, first, try, idtac, repeat},
    % Comments delimiters
    morecomment=[s]{(*}{*)},
    % Spaces are not displayed as a special character
    showstringspaces=false,
    % String delimiters
    morestring=[b]",
    morestring=[d]',
    % Size of tabulations
    tabsize=3,
    % Enables ASCII chars 128 to 255
    extendedchars=false,
    % Case sensitivity
    sensitive=true,
    % Automatic breaking of long lines
    breaklines=false,
    % Default style fors listings
    basicstyle=\small,
    % Position of captions is bottom
    captionpos=b,
    % flexible columns
    columns=[l]flexible,
    % Style for (listings') identifiers
    identifierstyle={\ttfamily\color{black}},
    % Style for declaration keywords
    keywordstyle=[1]{\ttfamily\color{dkviolet}},
    % Style for gallina keywords
    keywordstyle=[2]{\ttfamily\color{dkgreen}},
    % Style for sorts keywords
    keywordstyle=[3]{\ttfamily\color{ltblue}},
    % Style for tactics keywords
    keywordstyle=[4]{\ttfamily\color{dkblue}},
    % Style for terminators keywords
    keywordstyle=[5]{\ttfamily\color{dkred}},
    % Style for strings
    stringstyle=\ttfamily,
    % Style for comments
    commentstyle={\ttfamily\color{dkgreen}},
    literate=
    {\\forall}{{\color{dkgreen}{$\forall\;$}}}1
    {\\exists}{{$\exists\;$}}1
    {<-}{{$\leftarrow\;$}}1
    {=>}{{$\Rightarrow\;$}}1
    {==}{{\code{==}\;}}1
    {==>}{{\code{==>}\;}}1
    {->}{{$\rightarrow\;$}}1
    {<->}{{$\leftrightarrow\;$}}1
    {<==}{{$\leq\;$}}1
    {\#}{{$^\star$}}1 
    {\\o}{{$\circ\;$}}1 
    {\@}{{$\cdot$}}1 
    {\/\\}{{$\wedge\;$}}1
    {\\\/}{{$\vee\;$}}1
    {++}{{\code{++}}}1
    {~}{{$\sim$}}1
    {\@\@}{{$@$}}1
    {\\mapsto}{{$\mapsto\;$}}1
    {\\hline}{{\rule{\linewidth}{0.5pt}}}1
}[keywords,comments,strings]

% Define other code styles
\lstdefinestyle{mystyle}{
    commentstyle=\color{dkgreen},
    keywordstyle=\color{magenta},
    numberstyle=\tiny\color{codegray},
    stringstyle=\color{codepurple},
    basicstyle=\ttfamily\footnotesize,
    breakatwhitespace=false,         
    breaklines=true,                 
    captionpos=b,                    
    keepspaces=true,                 
    numbers=left,                    
    numbersep=5pt,                  
    showspaces=false,                
	    showstringspaces=false,
	    showtabs=false,                  
	    tabsize=2
	}

\definecolor{lightgrey}{RGB}{240,240,240}

\lstdefinelanguage{Lambdapi}
{
  inputencoding=utf8,
  extendedchars=true,
  numbers=none,
  numberstyle={},
  tabsize=2,
  basicstyle={\ttfamily\small\upshape},
  backgroundcolor=\color{lightgrey},
  keywords={abort,admit,admitted,apply,as,assert,assertnot,associative,assume,begin,builtin,change,commutative,compute,constant,debug,end,eval,fail,flag,focus,generalize,have,in,induction,inductive,infix,injective,left,let,notation,off,on,opaque,open,orelse,prefix,print,private,proofterm,protected,prover,prover_timeout,quantifier,refine,reflexivity,repeat,require,rewrite,right,rule,sequential,set,simplify,solve,symbol,symmetry,type,TYPE,unif_rule,verbose,why3,with},
  sensitive=true,
  keywordstyle=\color{blue},
  morecomment=[l]{//},
  morecomment=[n]{/*}{*/},
  commentstyle={\itshape\color{red}},
  string=[b]{"},
  stringstyle=\color{orange},
  showstringspaces=false,
  literate=
  {λ}{$\lambda$}1
  {↪}{$\hookrightarrow$}1
  {→}{$\rightarrow$}1
  {Π}{$\Pi$}1
  {≔}{$\coloneqq$}1
  {⊢}{$\vdash$}1
  {≡}{$\equiv$}1
  {𝔹}{$\mathbb{B}$}1
  {𝕃}{$\mathbb{L}$}1
  {ℕ}{$\mathbb{N}$}1
  {α}{$\alpha$}1
  {β}{$\beta$}1
  {η}{$\eta$}1
  {π}{$\pi$}1
  {τ}{$\tau$}1
  {ω}{$\omega$}1
  {∧}{$\wedge$}1
  {≤}{$\le$}1
  {≠}{$\neq$}1
  {∉}{$\notin$}1
  {×}{$\times$}1
  {⋅}{$\cdot$}1
  {⊑}{$\sqsubseteq$}1
}

% EBNF style definition
\lstdefinelanguage{EBNF}{
  morekeywords={::=},
  sensitive=true,
  morecomment=[s]{(*}{*)},
  morestring=[b]"
}

\definecolor{comment}{RGB}{0,128,0}
\definecolor{keyword}{RGB}{0,0,255}
\definecolor{nonterminal}{RGB}{128,0,128}
\definecolor{terminal}{RGB}{165,42,42}

\lstdefinestyle{ebnfstyle}{
  language=EBNF,
  basicstyle=\ttfamily,
  commentstyle=\color{comment}\itshape,
  keywordstyle=\color{keyword}\bfseries,
  stringstyle=\color{terminal},
  identifierstyle=\color{nonterminal},
  frame=single,
  breaklines=true,
  showstringspaces=false,
  tabsize=2
}

\newcommand{\bnf}[1]{\textsf{<}#1\textsf{>}}
\lstset{style=mystyle}

% Title and author information
\title{First-Class Constructor Subsets for Pattern Matching on Indexed Families}

\author{Tiago Campos}{UFMG, Brazil}{camposferreiratiago@gmail.com}{https://orcid.org/0000-0002-1825-0097}{}

\authorrunning{T. Campos}

\Copyright{Tiago Campos}

\ccsdesc[300]{Theory of computation~Program verification}
\ccsdesc[500]{Theory of computation~Type theory}

\keywords{Datatype, Type Constructors, Constructor Subtyping}

\begin{document}

\maketitle
\bibliographystyle{plainurl}% the mandatory bibstyle

\begin{abstract}
Pattern matching is a powerful feature of functional and dependently typed languages, yet programmers often encounter unreachable clauses. This either leads to unsafe partial functions or, in safe settings, to proof obligations that require tedious resolution of trivial cases. The problem is particularly acute when programs operate on subsets of an inductive datatype. Standard encodings pair data with propositions or refactor datatypes into indexed families, but both approaches complicate composition and proof automation.

We present a subtyping discipline for constructor subsets, where first-class signatures describe which constructors of a family are available. These signatures make the intended subset explicit at the type level, enabling pattern matching to rule out impossible branches without generating spurious obligations. Our prototype indicates that this approach can eliminate many trivial proof obligations while preserving type safety.

We implemented a small proof assistant that supports constructor subsets (signatures), and we developed a small prototype that translates its theory to the Lambdapi framework. We can encode the entire type system in fewer than 600 lines, supporting our claim that the approach is simple and lightweight.

%Edit: Clarified (in the Introduction below) what ``available'' means operationally: signatures constrain the \emph{outermost} constructor observed by pattern matching, which is exactly what removes unreachable branches without forcing deep/recursive constructor closure.
\end{abstract}

\section{Introduction}
Subtyping is a fundamental concept across programming languages that enables the specialization of type definitions through hierarchical relationships between supertypes and subtypes. At a high level, a subtype describes a ``smaller'' collection of values than its supertype. While term-level subtyping is extensively implemented in many languages, subtyping between constructors of inductive types remains relatively unexplored, especially in the presence of dependent types.

A recurring difficulty arises when an inductive datatype is only partially available in a given context. For example, a function may be intended to consume lists that are known to be non-empty, or vectors whose length is positive. In most dependently typed languages the programmer faces a choice: either write an unsafe partial function, or thread propositions through the program and discharge proof obligations stating that certain constructors are impossible. The latter approach leads to unreachable clauses in pattern matches and to proofs that are often routine but cumbersome.

One common technique is to maintain the original datatype together with a proposition that rules out some constructors. Another approach, illustrated in Figures~\ref{fig:coq-obligation} and~\ref{fig:coq-indexed}, uses indexed datatypes to restrict which constructors may appear. Both techniques have well-known drawbacks. The first ties ordinary programs to logical predicates, complicating function composition and proof reuse. The second introduces additional indices and can interact badly with proof irrelevance and unification; it may also hinder the detection of unreachable branches when constructor information flows through computations.

Constructor subtyping has been explored in other settings. In O'Haskell~\cite{Nordlander987777}, records and algebraic datatypes can be equipped with an optional subtyping relation inferred from the usage of constructors and fields. For example, a non-empty list type can be extended with an empty constructor to obtain a standard list type. More recent work on zero-cost constructor subtyping in Cedille~\cite{MarmadukeJenkinsStump,jenkins2019elaborating}, on order-sorted inductive types and overloading~\cite{BatheGillesFradeMaria,barthe}, and on algebraic subtyping for extensible records~\cite{MarquesFloridoVasconcelos2024} shows how rich subtyping disciplines can be added to strongly typed systems. However, these systems are tailored to specific core calculi or to non-dependent settings, and do not directly address the combination of constructor subtyping with indexed families in a general dependently typed setting.

%Edit: Clarified the intended meaning of constructor subsets to match the operational goal (coverage of pattern matching) and the later examples.
In this paper we focus on a simple but expressive fragment of constructor subtyping aimed at eliminating trivial proof obligations. Rather than allowing arbitrary constructor overloading, we work with \emph{constructor subsets}: first-class types of the form $\{T :: \Phi\}$, where $T$ is an inductive family and $\Phi$ is a finite list of its constructors. Intuitively, $\{T :: \Phi\}$ denotes those values of $T$ whose \emph{outermost (head) constructor}, when inspected in canonical form, is drawn from $\Phi$. This is the information that coverage checking for pattern matching consumes directly, and it is exactly what makes types such as ``non-empty lists'' inhabited (the tail may still be empty). 

Our contributions are as follows.
\begin{itemize}
\item We define a subtyping discipline for constructor subsets as a conservative extension of a $\lambda\Pi$-style dependent type theory with inductive families. The system includes typing and subtyping rules for constructor signatures and a dedicated elimination rule for pattern matching on signature types.
\item We show, through a series of examples, that constructor subsets capture common programming idioms such as non-empty lists and simple invariants on inductive families, while avoiding the proof obligations that arise when these invariants are expressed with explicit propositions.
\item We describe a two-layer prototype: a Lambdapi backend encoding of constructor signatures and a Haskell transpiler that translates a lightweight surface fragment into signature-aware LambdaPi terms. This architecture keeps programs concise while preserving a direct correspondence with Rules~1--7.
\end{itemize}

\paragraph*{Structure of the paper.}
Section~\ref{sec:rules} introduces the core type system and subtyping rules.
Section~\ref{sec:cases} presents case studies illustrating how constructor subsets simplify programming with indexed families.
Section~\ref{sec:implementation} sketches the implementation and discusses practical considerations, and Section~\ref{sec:conclusion} concludes with directions for future work.

\begin{figure}
\begin{lstlisting}[language=Coq]
Inductive list (A : Set) : Set :=
    |cons : A -> list A -> list A
    |empty : list A.

Definition head {A} (x : list A) : 
    x <> empty A -> A := 
    match x return x <> empty A -> A with
        |cons _ h l => fun _ => h
        |empty _ => 
            fun f => 
                match (f eq_refl) with end
end.
\end{lstlisting}
\caption{Constructor exclusion in Coq with a proof obligation.}
\label{fig:coq-obligation}
\end{figure}

\begin{figure}
\begin{lstlisting}[language=Coq]
Inductive list' (A : Set) : bool -> Set :=
    |cons' : forall x, A -> 
        list' A x -> list' A true
    |empty' : list' A false.

Definition head' {A} (x : list' A true) : A := 
    match x with
       |cons' _ _ h l => h
   end.
\end{lstlisting}
\caption{Constructor exclusion in Coq using an indexed datatype.}
\label{fig:coq-indexed}
\end{figure}

Many type systems are inspired by constructor subtyping to address constructor overloading, allowing the same constructor name to be shared by different datatypes, even in pattern matching. However, its naive use can lead to problems with subject reduction, as observed by Frade~\cite{mariaJoaoFrade}. Subject reduction is the property that during the normalization of a term its type is preserved; overloading constructors without care may violate this property. For the sake of simplicity, in this work we do not consider overloading between distinct datatypes. Instead we restrict subtyping to relations induced by subsets of constructors of a single inductive family. We believe this fragment is already expressive enough to avoid the trivial proof obligations discussed above, while keeping the meta-theory and implementation comparatively simple.

\section{Type Rules}
\label{sec:rules}
In this section, we introduce a simplified type system, enhanced with first-class signature subtyping support. This exposition assumes a basic familiarity with the concepts of type theory, especially dependent type theory. We use a standard typing judgment $\Gamma \vdash t : A$ for terms and a subtyping judgment $\Gamma \vdash A \sqsubseteq B$ between types. Contexts $\Gamma$ map variables to types, and $Type$ denotes the universe of small types.

The subtyping relation is inspired by the subsumption rule of Aspinall and Compagnoni~\cite{AspinallCompagnoni}:
\begin{center}
\begin{mathpar}
\inferrule* []{\Gamma \vdash t : A \quad \Gamma \vdash A \sqsubseteq B}{\Gamma \vdash t : B}
\end{mathpar}
\end{center}

We now fix some notation for telescopes and constructor environments. A \emph{telescope} $\Delta$ is a sequence of typed variables
\[
  \Delta = (a_{1} : A_{1}) \dots (a_{n} : A_{n}).
\]
Given such a telescope and a type family $T$, we write $T\,\Delta$ for the application $T\,a_{1}\,\dots\,a_{n}$. When a telescope is used only to record the indices of a datatype we write $\Delta^{*}$; in examples we use $Vector\,A\,n$ with $\Delta = (A : Type)(n : Nat)$ and $\Delta^{*} = (A,n)$~\cite{Norell2007TowardsAP}.

To refer to individual arguments in a telescope we use subscripts: if $\Delta = (a_{1} : A_{1})\dots(a_{n} : A_{n})$ then $\Delta_i$ denotes $(a_i : A_i)$ and $\Delta^{*}_{i}$ denotes the corresponding index, for $1 \leq i \leq n$. For a term $T$ we use the side condition $T_{\beta\eta}$ to indicate that we inspect the $\beta\eta$-normal form of $T$.

We write $\mathcal{C}_{\text{all}}$ for the finite set of static constructor declarations provided by the programmer. 
A static definition stores a constructor name. We do not rely on any particular labeling scheme; we only need to decide constructor equality, so one may use unique names or hashes.
The order of static definitions is stored in $\mathcal{C}_{\text{all}}$; in particular, mutually recursive definitions can be handled by first registering every constructor in $\mathcal{C}_{\text{all}}$ and only then checking the rules below. A \emph{signature} $\Phi$ is a finite list of constructor names drawn from $\mathcal{C}_{\text{all}}$, and a type of the form $\{T\,\Delta^{*} :: \Phi\}$ maps each constructor in $\Phi$ to the restricted type $T\,\Delta^{*}$. Formation is intentionally liberal: signatures may include constructors that are incompatible with the current index instance $\Delta^{*}$, which yields uninhabited refinements but does not introduce inconsistency. We use a side predicate $\mathsf{AGAINST}$ only for elimination-time pruning (Rule~7), not for signature formation. Because $\mathcal{C}_{\text{all}}$ is finite, the side conditions that test membership in $\mathcal{C}_{\text{all}}$ are decidable.

%Edit: The rules below were rewritten to make the quantification over constructors and their result indices explicit (the intended meaning of the original rules), while preserving the same information/content.
\begin{center}
\begin{mathpar}

\inferrule* [left=Rule 1]{%
\Gamma \vdash T\:\Delta^{*} : Type \quad \Phi = (C_{1}, \ldots, C_{n}) \quad (C_{1} \ldots C_{n}) \subseteq \mathcal{C}_{\text{all}} \\
\forall i,\, 1 \le i \le n,\ \Gamma \vdash C_{i} : \Delta_{i} \rightarrow T\:\Delta_{i}'^{*}%
}{%
\Gamma \vdash \{T \Delta^{*} :: \Phi \} : Type%
} \\

\inferrule* [left=Rule 2]{%
\Gamma \vdash T\:\Delta^{*} : Type%
}{%
\Gamma \vdash \{T\:\Delta^{*} :: \Phi \} \sqsubseteq T\:\Delta^{*}%
} \\

%Edit: Rule 3 was rewritten for clarity: it introduces a constructor application into the corresponding signature type when the constructor is listed in the signature and the result indices match (up to $\beta\eta$) the signature's indices.
\inferrule* [left=Rule 3, right=$C\:\in\:\mathcal{C}_{\text{all}}$]{%
\Gamma \vdash T\:\Delta^{*} : Type \\
\Gamma\:\vdash\:C:\Delta_{C} \rightarrow T\:\Delta_{C}'^{*} \quad C \in \Phi \\
\Gamma \vdash \overrightarrow{u} : \Delta_C \\
(T\:\Delta_{C}'^{*}[\overrightarrow{u}/\Delta_C])_{\beta\eta} = (T\:\Delta^{*})_{\beta\eta}%
}{%
\Gamma \vdash C \:\overrightarrow{u} : \{T\:\Delta^{*} :: \Phi \}%
} \\

%Edit: Rule 4 kept the original intent but makes the index application of $T$/$T'$ explicit.
\inferrule* [left=Rule 4]{%
(T\:\Delta^{*})_{\beta\eta} = (T'\:\Delta^{*})_{\beta\eta} \quad \Phi' \subseteq \Phi%
}{%
\Gamma \vdash \{T' \Delta^{*} :: \Phi' \} \sqsubseteq \{T \Delta^{*} :: \Phi \}%
} \\

\inferrule* [left=Rule 6]{\Gamma \vdash A' \sqsubseteq A \\ \Gamma, x:A' \vdash B \sqsubseteq B'}{\Gamma \vdash (x : A) \rightarrow B \sqsubseteq (x : A') \rightarrow B'}

\end{mathpar}
\end{center}

\paragraph*{Rule 5}
We keep the following rule as an implementation-level checking convenience (not part of the declarative calculus above). It is derivable from standard application typing plus subsumption, but it is convenient to have it as a direct check in the implementation to avoid unnecessary normalization if $F$ is not a constructor application:
\begin{mathpar}
\inferrule* [left=Rule 5, right=$F\notin\mathcal{C}_{\text{all}}$]{%
\Gamma \vdash F : \Delta \rightarrow A \\
\Gamma \vdash \overrightarrow{u} : \Delta \\
\Gamma \vdash A[\overrightarrow{u}/\Delta] \sqsubseteq \{T\:\Delta^{*} :: \Phi\}%
}{%
\Gamma \vdash F\:\overrightarrow{u} : \{T\:\Delta^{*} :: \Phi\}%
}
\end{mathpar}

A term annotated as $T_{\beta\eta}$ is considered in $\beta\eta$-normal form. In practice we only normalize enough to inspect the outermost constructor when applying Rule~3 (and the algorithmic Rule~5 check in the implementation); there is no global requirement that all terms be strongly normalizing. The choice of how aggressively to normalize is left to the implementation and presents the usual trade-off between completeness of simplification and performance.

It is straightforward to see that $\text{cons}\:\Delta$ has type $\{\text{Vector}\:\Delta^* :: |\text{cons}\}$. One may either infer this type and then check it against an expected type $T$, or attempt to check $\text{cons}\:\Delta$ directly against $T$. In our prototype we follow the former strategy, inferring first and checking second, because it tends to reduce the number of normalization steps needed during type checking.

%Edit: Added a short operational reading to align the rules with the coverage-checking motivation and the head-constructor interpretation stated in the Introduction.
\paragraph*{Operational reading.}
Under the usual canonical forms lemma for inductive values, a closed value of a signature type $\{T\:\Delta^{*} :: \Phi\}$ reduces (in weak head normal form) to a constructor application $C\:\overrightarrow{u}$ with $C \in \Phi$ and $(T\:\Delta_{C}'^{*}[\overrightarrow{u}/\Delta_C])_{\beta\eta} = (T\:\Delta^{*})_{\beta\eta}$. Thus, signature types refine \emph{which head constructors} can appear at a given program point, exactly the information required to remove unreachable branches in pattern matching.

Signatures may contain constructors that are incompatible with the current indices. For example,
\[
T \coloneqq \{\text{Vector}\:A\:(S\;N)\:::\:|\text{empty}|\text{cons}\}
\]
is well-formed, but $\text{empty}$ cannot inhabit $T$. These trivial subsets are harmless: they are simply uninhabited refinements, and there is no rule that introduces $\{T\:\Delta^{*}::\Phi\}$ from $T\:\Delta^{*}$ without constructor evidence.

To eliminate unreachable branches in pattern matching, we define a pruned constructor list:
\begin{math}
\Phi_{\mathsf{ok}}
\;\coloneqq\;
\mathsf{OK}(\Gamma,T,\Delta^{*},\Phi)
\;\equiv\; 
\bigl[\; C \in \Phi \ \big|\ 
  C \in \mathcal{C}_{\text{all}}
  \ \land\ 
  \Gamma \vdash C : \Delta_C \rightarrow T\,\Delta_C'^{*}
  \ \land\ \\
  \forall j \in \{1,\dots,|\Delta^{*}|\}.\;
    \mathsf{AGAINST}\bigl((\Delta^{*}_j),(\Delta_{C,j}'^{*})\bigr)
\;\bigr]
\end{math}, where $\Delta^{*}$ denotes the index tuple of type $T$.

The predicate $\mathsf{AGAINST}$ is given below:
\[
\mathsf{AGAINST}(p,t) \text{ is defined by the following decision procedure:}
\]
\[
\mathsf{AGAINST}(p,t) \;\equiv\;
\begin{cases}
\top
  & \text{if } p \text{ is a variable} \\
\bigwedge_{i=1}^{k}\mathsf{AGAINST}(p_i,t_i)
  & \text{if } p = c\,p_1\ldots p_k \text{ and } t = c\,t_1\ldots t_k \\
\bot
  & \text{if } p = c\,\overrightarrow{p'} \text{ and } t = c'\,\overrightarrow{t'}
    \text{ with } c \neq c' \\
\top
  & \text{otherwise.}
\end{cases}
\]

Rule~2 allows us to forget signature information at the same index instance: any value of type $\{T\:\Delta^{*} :: \Phi' \}$ can be viewed as a value of the underlying type $T\:\Delta^{*}$. In this sense constructor subsets behave like \emph{phantom types}~\cite{fluet2005practical}: the extra information carried by the signature is present at type-checking time but is not reflected in the runtime representation of values.
	
Although phantom types are not widely known across programming languages and are used primarily in Haskell, they have important properties for representing types that carry no runtime information; that is, they are purely erased types~\cite{fluet2005practical}. In our system, there is no dedicated elimination rule for pure phantom signatures, and Rule~7 only permits case analysis on signatures that still expose constructors.

A constructor $C$ must adhere to certain restrictions in its definition. For example, the signature (i.e., a type $\{T\:\Delta^* \:::\:\Phi\}$) cannot occur freely in $C$. This restriction prohibits constructor types such as $\text{succ}:\{\text{nat} :: |\text{succ}\:|\:0\} \rightarrow \{\text{nat} :: |\text{succ}\:|\:0\}$, which would lead to problematic self-reference, as $\text{succ}$ would require itself in its own definition.

It is important to note that even $succ: nat \rightarrow nat$ constitutes an invalid definition in this system. This is because its predecessor argument is a phantom type; therefore, for $C$ to satisfy the inductive criteria, a signature must occur freely only in a positive position in the $\Delta$ of $\text{succ}$.

%Edit: Added a clarifying remark about what is (and is not) covered regarding positivity/inductive well-formedness, without removing the original discussion.
\paragraph*{Remark (inductive well-formedness).}
The previous paragraph sketches the kind of restrictions that an implementation must enforce so that inductive definitions remain well-founded (e.g.\ positivity/strictness conditions and well-scoped use of signatures in recursive arguments). Our prototype enforces a conservative check in the spirit of standard positivity conditions; however, we do not attempt to present a complete account of inductive well-formedness for signatures in this paper.

We do not cover index matching (from dependent (co)pattern matching) rules and conversion details in this work, as they are beyond our current focus. We leave the implementation of these principles to the authors' discretion.

\begin{center}
\begin{mathpar}
\inferrule* [left=Rule 7]{%
\Gamma \vdash T\:\Delta^* : Type
\quad
\Gamma \vdash Q : Type
\\
\Gamma \vdash M :\{T\:\Delta^* :: \Phi\}
\quad
\Phi_{\mathsf{ok}} = \mathsf{OK}(\Gamma,T,\Delta^{*},\Phi)
\\
\Phi_{\mathsf{ok}} \subseteq \Psi \subseteq \Phi
\quad
\Psi = (D_{1},\ldots,D_{m})
\\
\forall k,\ 1 \le k \le m,\ 
\Gamma \vdash D_{k} : \Delta_{D_k} \rightarrow T\:\Delta_{D_k}'^*
\\
\forall k,\ 1 \le k \le m,\ 
\Gamma \vdash N_{k} : \Delta_{D_k} \rightarrow Q%
}{%
\Gamma \vdash
\textsf{case}\; M\; \textsf{of}\; Q\;
\{\, D_{k}\,\Delta_{D_k} \Rightarrow N_{k}\,\Delta_{D_k} \,\}_{k=1}^{m}
: Q%
} \\
\end{mathpar}
\end{center}

Rule~7 requests branches only for constructors in a branch list $\Psi$ that covers the obviously compatible constructors: $\Phi_{\mathsf{ok}} \subseteq \Psi \subseteq \Phi$. In practice one chooses $\Psi = \Phi_{\mathsf{ok}}$, which removes unreachable branches while keeping the rule conservative.
There is yet a fancy way of dealing with signatures that contains trivial subsets, we could normalize them to their pruned version $\Phi_{\mathsf{ok}}$.
We model this as an optional normalization step $\longrightarrow_{\phi}$ on signature types.

Any conservative $\mathsf{AGAINST}$ definition that is sound for pruning is admissible: 
if it rejects a constructor at some index position, then Rule~3's index equality cannot hold for that constructor at those indices.

\begin{mathpar}
\inferrule*[left=\textsc{Norm}$_\phi$]{
  \Phi_{\mathsf{ok}} = \mathsf{OK}(\Gamma,T,\Delta^{*},\Phi)
}{
  \{T\,\Delta^{*} :: \Phi\} \longrightarrow_{\phi} \{T\,\Delta^{*} :: \Phi_{\mathsf{ok}}\}
}
\end{mathpar}

As a further (optional) convenience, one may treat $\phi$-normalization as part of
definitional equality. We write $\Gamma \vdash A \equiv_{\beta\eta\phi} B$ for
conversion that closes $\beta\eta$-conversion with $\phi$-simplification on signatures.
This can be read operationally as: normalize both sides by $\beta\eta$ and repeatedly
apply \textsc{Norm}$_\phi$ at signature nodes, then compare up to $\alpha$-equivalence.
We do not develop the meta-theory of $\equiv_{\beta\eta\phi}$ in this paper; although still simple to adhere to. In fact, many rules will be just simplified.
It is presented as a design option that matches the implementation strategy of keeping signatures pruned. Although there is a cost of augmenting the complexity of the formalism, since you need to carry a new conversion relation.

\begin{mathpar}
\inferrule*[left=Eq$_\phi$]{
  \Phi_{\mathsf{ok}} = \mathsf{OK}(\Gamma,T,\Delta^{*},\Phi)
}{
  \Gamma \vdash \{T\,\Delta^{*} :: \Phi\} \equiv_{\beta\eta\phi} \{T\,\Delta^{*} :: \Phi_{\mathsf{ok}}\}
}
\end{mathpar}

\begin{mathpar}
\inferrule* [left=Conversion$_\phi$]
{\Gamma \vdash t : A \\
 \Gamma \vdash B : s \\
 \Gamma \vdash A \equiv_{\beta\eta\phi} B}
{\Gamma \vdash t : B}
\end{mathpar}

$\equiv_{\beta\eta\phi}$ is the smallest congruence containing $\beta\eta$-conversion modulo $\mathrm{Eq}_{\phi}$. We state some properties of the rules 1-7 the system above:

%Edit: The statement below is preserved but clarified: the ``$\simeq$'' correspondence is best understood as an external encoding/translation into a host theory with $\Sigma$-types; the reverse direction is not derivable in the presented core rules without an additional refinement-introduction principle.
\begin{theorem}[Head-constructor refinement]
Let $T$ be an inductive family with constructor set $C$, and let $\Phi \subseteq C$ be a signature at indices $\Delta^{*}$.
Assume
\[
  \Gamma \vdash \{T\:\Delta^{*}::\Phi\} : Type
  \quad\text{and}\quad
  \Phi = \Phi_{\mathsf{ok}}.
\]
Define $D : T\:\Delta^{*} \to Type$ by inspecting only the \emph{head constructor} of
the canonical form: for each constructor application $C'\,\overrightarrow{u}$,
\[
  D(C'\,\overrightarrow{u}) \;\equiv\;
  \begin{cases}
    \top & \text{if } C' \in \Phi,\\
    \bot & \text{if } C' \notin \Phi.
  \end{cases}
\]
Then there is a definable map
\[
  f : \{T\:\Delta^{*}::\Phi\} \to \Sigma(x:T\:\Delta^{*}), D(x).
\]
Moreover, in any host theory that validates refinement introduction
(from $x:T\:\Delta^{*}$ and $D(x)$ to $x:\{T\:\Delta^{*}::\Phi\}$), this map is part of an
isomorphism $\{T\:\Delta^{*}::\Phi\} \simeq \Sigma(x:T\:\Delta^{*}), D(x)$.
\end{theorem}
\begin{proof}
Define $f(x) \coloneqq (x, \star)$.
To justify $\star : D(x)$, use the canonical-forms property for signature
types: if $x : \{T\:\Delta^{*}::\Phi\}$ is a value, its canonical form is
$C\,\overrightarrow{u}$ with $C \in \Phi$, hence $D(C\,\overrightarrow{u}) \equiv \top$.

For the converse direction, assume refinement introduction:
from $(y,d) : \Sigma(x:T\:\Delta^{*}), D(x)$ we obtain $y : \{T\:\Delta^{*}::\Phi\}$ and define
$g(y,d) \coloneqq y$. Then $f$ and $g$ are mutually inverse.
\end{proof}

\begin{mathpar}
\inferrule* [left=$\xi$-$\cdot$app$_1$]
  {L \longrightarrow L'}
  {L \cdot M \longrightarrow L' \cdot M}

\inferrule* [left=$\xi$-$\cdot$app$_2$]
  {M \longrightarrow M'}
  {V \cdot M \longrightarrow V \cdot M'}

\inferrule* [left=$\xi$-$\beta$]
  { }
  {(\lambda x \Rightarrow N) \cdot V \longrightarrow N [ x := V ]}
\end{mathpar}

\begin{mathpar}
\inferrule* [left=$\xi$-$case$]
  {v \longrightarrow v'}
  {case\:v\:of\:Q\:\{C_{i}\:\Delta_{i} \Rightarrow N_{i}\:\Delta_{i} {,}...\} \longrightarrow case\:v'\:of\:Q\:\{C_{i}\:\Delta_{i} \Rightarrow N_{i}\:\Delta_{i} {,}...\} }

\inferrule* [left=$\xi$-$case'$]
  {\quad}
  {case\:(C_{i}\:\Delta' )\:of\:Q\:\{C_{i}\:\Delta_{i} \Rightarrow N_{i}\:\Delta_{i} {,}...\} \longrightarrow N_i\:\Delta'}
\end{mathpar}

Metavariables \(M,N,S\) range over terms; \(V,W\) range over values (lambda abstractions or constructor applications in WHNF).
In the rule \(\xi\)-case above, the symbols \(v,v'\) denote arbitrary terms (not values).

\begin{theorem}[Progress]
If $\cdot \vdash M : \{T\:\Delta^{*} :: \Phi\}$, then either $M$ is a value or there exists $M'$ such that $M \longrightarrow M'$.
\end{theorem}
\begin{proof}
We reason by structural induction on $M$ and case analysis on its form.
\begin{itemize}
\item \emph{Application $M = M_1 \cdot M_2$.} If either $M_1$ or $M_2$ is not a value, the induction hypothesis yields $M_i'$ with $M_i \longrightarrow M_i'$, and we use the corresponding congruence rule $\xi$-app$_i$ to obtain a step for $M$. Otherwise both $M_1$ and $M_2$ are values. If $M_1$ is a lambda abstraction $\lambda x. N$, we can perform a $\beta$-reduction step. If $M_1$ is a constructor, then $M$ is already a value (an applied constructor) and the conclusion holds.
\item \emph{Case expression $M = \text{case}\;v\;\text{of}\;Q\;\{C_i\:\Delta_i \Rightarrow N_i\:\Delta_i, \ldots\}$.} If $v$ is not a value, the induction hypothesis gives $v'$ with $v \longrightarrow v'$, and we use rule $\xi$-case. If $v$ is a value, typing gives $v : \{T\:\Delta^{*} :: \Phi\}$, so by canonical forms $v = C\,\overrightarrow{u}$ with $C \in \Phi$ and matching result indices. Therefore $C$ is not rejected by $\mathsf{AGAINST}$, hence $C \in \Phi_{\mathsf{ok}}$. Rule~7 requires a branch list $\Psi$ such that $\Phi_{\mathsf{ok}} \subseteq \Psi$, so the branch for $C$ is present and we apply \(\xi\)-case'.
\item \emph{Other forms.} For variables, constructors, and abstractions the canonical forms analysis from the typing rules shows that well-typed closed terms of signature type are either values or reduce by one of the cases above.
\end{itemize}
In all cases, a closed term of type $\{T\:\Delta^{*} :: \Phi\}$ is either a value or can take a reduction step.
\end{proof}

\begin{theorem}[Preservation]
If $\Gamma \vdash M : R$ and $M \longrightarrow M'$, then $\Gamma \vdash M' : R$.
\end{theorem}
\begin{proof}
We proceed by induction on the evaluation derivation, considering the last reduction rule used.
\begin{itemize}
\item \emph{$\beta$-reduction.} If $M = (\lambda x. N) \cdot V$ and $M' = N[x := V]$, typing gives $\Gamma \vdash \lambda x. N : (x : A) \rightarrow B$ and $\Gamma \vdash V : A$. By the substitution lemma we obtain $\Gamma \vdash N[x := V] : B[x := V]$, which is the same type as $M$.
\item \emph{$\xi$-app$_1$ and $\xi$-app$_2$.} In these cases we reduce a proper subterm of $M$. The induction hypothesis states that the reduced subterm preserves its type; reapplying the application typing rule reconstructs a typing derivation for $M'$ with the same type $R$.
\item \emph{$\xi$-case.} Here $M = \text{case}\;v\;\text{of}\;Q\;\{\dots\}$ and $M' = \text{case}\;v'\;\text{of}\;Q\;\{\dots\}$ with $v \longrightarrow v'$. By the induction hypothesis $\Gamma \vdash v' : \{T\:\Delta^{*} :: \Phi\}$, and Rule~7 then re-establishes $\Gamma \vdash M' : Q = R$.
\item \emph{$\xi$-case'.} In this case $M = \text{case}\;(C_i\:\Delta')\;\text{of}\;Q\;\{\dots\}$ reduces to $M' = N_i\:\Delta'$. Rule~7 gives $\Phi_{\mathsf{ok}} \subseteq \Psi \subseteq \Phi$ for the branch list and typing premises $\Gamma \vdash N_k : \Delta_{D_k} \rightarrow Q$ for every $D_k \in \Psi$. By canonical forms on the scrutinee, the head constructor satisfies $C_i \in \Phi_{\mathsf{ok}}$, hence $C_i \in \Psi$ and the corresponding branch exists. Instantiating with the actual arguments $\Delta'$ yields $\Gamma \vdash N_i\:\Delta' : Q$, so $M'$ has the same type as $M$.


\end{itemize}
All other reduction rules follow the same pattern: either a proper subterm reduces and the induction hypothesis applies, or a local computation such as $\beta$-reduction is justified by the typing rules. In each case the result has the same type $R$ as the original term.
\end{proof}

\section{Case Examples}
\label{sec:cases}
We introduce a minimal language and its implementation to demonstrate the subtyping system. The abstract syntax of terms is sketched in Figure~\ref{fig:term-syntax}; in the examples below we use a lightweight ML-style surface syntax.
	
\begin{figure}[!ht]
\begin{align*}
\text{Term} &::= \text{Var}(x) \mid \text{App}(f, a) \mid \text{Lam}(x, e) \mid \text{Pi}(x, A, B) \\
&\mid \text{Constr}(T, \overrightarrow{c}) \mid \text{Match}(e, T, \overrightarrow{p \to e'}) \mid \text{Notation}(e, T)
\end{align*}
	
\caption{Term syntax of the example language}
\label{fig:term-syntax}
\end{figure}

\paragraph{Mini-language overview.}
We built a proof assistant in Haskell and implemented the type system described above, along with other features such as unification.

Programs consist of a sequence of \emph{static declarations} followed by regular statements. Static declarations register constructors in $\mathcal{C}_{\text{all}}$ (Rule~1), whereas dynamic statements bind names to terms typed by the rules from Section~\ref{sec:rules}. Terms use the usual $\lambda$-abstractions, function types, applications, and pattern matches; the only novelty is that constructor signatures $\{T :: \Phi\}$ appear explicitly in types. These informal descriptions suffice for the examples below; the full concrete syntax is implemented in the prototype but omitted here.

We exploit the capability to operate with dynamically defined constructor subsets without requiring trivial proof obligations. Since datatype signatures are first-class inhabitants in our system, we can instantiate varying constructor sets through definition aliasing. The following declarations register constructors for lists and expose two different signatures:
\begin{lstlisting}[
caption={Definition of empty and non-empty lists using first-class signatures},
language=ML,
basicstyle=\ttfamily\small,
keywordstyle=\bfseries\color{blue},
commentstyle=\itshape\color{gray},
stringstyle=\color{red},
numbers=left,
numberstyle=\tiny\color{gray},
frame=single
]
Static list : *> *.
Static empty : (A : *) -> (list A).
Static new : (A : *) -> A> {(list A) :: |new |empty}> (list A).
List |A :: *> * => {(list A) :: |new |empty}.
NonEmpty |A :: *> * => {(list A) :: |new}.
\end{lstlisting}
Rule~3 ensures that each constructor inhabits the declared signature, while Rule~4 provides $\textsf{NonEmpty}\:A \sqsubseteq \textsf{List}\:A$ because the latter merely adds the \texttt{empty} constructor. This formulation enables the definition of functions that operate on both general lists and non-empty lists:
\begin{lstlisting}[
caption={Length function compatible with both list variants},
language=ML,
basicstyle=\ttfamily\small,
keywordstyle=\bfseries\color{blue},
commentstyle=\itshape\color{gray},
stringstyle=\color{red},
numbers=left,
numberstyle=\tiny\color{gray},
frame=single
]
length |A ls :: (A : *) -> (List A)> Nat =>
[ls of Nat
|(empty _) => 0
|(new A head tail) => (+1 (length A tail))
].
\end{lstlisting}
Alternatively, we can define functions that exclusively accept non-empty lists, directly mirroring the running example:
\begin{lstlisting}[
caption={Subtyping application to eliminate trivial proof obligations},
language=ML,
basicstyle=\ttfamily\small,
keywordstyle=\bfseries\color{blue},
commentstyle=\itshape\color{gray},
stringstyle=\color{red},
numbers=left,
numberstyle=\tiny\color{gray},
frame=single
]
last |A ls :: (A : *) -> (NonEmpty A)> A =>
[ls of A
    |(new A head tail) => [tail of A
        |(empty _) => head
        |(new A head2 tail2) => (last A (new A head2 tail2))
    ]
].
insertsort |xs v :: (List Nat)> Nat> (NonEmpty Nat) =>
[xs of (NonEmpty Nat)
    |(empty _) => (new Nat v (empty Nat))
    |(new _ head tail) => [(gte head v) of (NonEmpty Nat)
        |false => (new Nat head (insertsort tail v))
        |true => (new Nat v (new Nat head tail))
    ]
].
def list_inserted_has_a_last_element |xs v :: (List Nat)> Nat> Nat =>
(last Nat (insertsort xs v))
\end{lstlisting}
Rules~2 and~7 explain why \texttt{last} needs only one clause: the type of \texttt{ls} is $\{(List\:A) :: |new\}$, so no \texttt{empty} branch is requested. The helper \texttt{list_inserted_has_a_last_element} demonstrates that the sorted list inherits the non-empty signature with no auxiliary proofs---exactly the two obligations eliminated relative to Figure~\ref{fig:coq-obligation}.
	
One might assume the following representation is correct, as previously discussed, yet Rule~2 warns against using phantom predecessors:
	
\begin{lstlisting}[
caption={Attempt with phantom predecessor},
language=ML,
basicstyle=\ttfamily\small,
keywordstyle=\bfseries\color{blue},
commentstyle=\itshape\color{gray},
stringstyle=\color{red},
numbers=left,
numberstyle=\tiny\color{gray},
frame=single
]
Static nat : *.
Static 0 : nat.
Static +1 : nat> nat.
Nat :: {nat :: |0 |+1}.
\end{lstlisting}

Here the predecessor of \texttt{+1} would live in $\{nat :: |0 | +1\}$ but the constructor type mentions only the phantom \texttt{nat}. The checker's algorithmic Rule~5 rejects this definition, forcing us to use the recursive variant below:

\begin{lstlisting}[
caption={Recursive definition accepted by the algorithmic Rule 5 check},
language=ML,
basicstyle=\ttfamily\small,
keywordstyle=\bfseries\color{blue},
commentstyle=\itshape\color{gray},
stringstyle=\color{red},
numbers=left,
numberstyle=\tiny\color{gray},
frame=single
]
Static nat : *.
Static 0 : nat.
Static +1 : Nat> nat.
Nat :: {nat :: |0 |+1}.
\end{lstlisting}
Because the predecessor now mentions the signature \texttt{Nat}, Rule~3 can derive the constructor typing judgment directly. The pruning function used by Rule~7 then removes obviously incompatible branches when matching on such signatures, avoiding unnecessary proof obligations~\cite{martin2021intuitionistic}.

\section{Implementation}
\label{sec:implementation}
Our implementation now has three components.
The first component is the minimal language implemented above, which supports inductive datatypes in a superset of the $\lambda \Pi$ calculus modulo (Figure~\ref{fig:lambdapi-figure3}).

\begin{figure}[h]
\begin{center}
\begin{mathpar}
\inferrule* [left=Sort]{ }{\Gamma;\Delta \vdash Type : Kind}

\inferrule* [left=Variable]{(x : A) \in \Delta}{\Gamma;\Delta \vdash x : A}

\inferrule* [left=Constant]{(c : A) \in \Gamma}{\Gamma;\Delta \vdash c : A}

\inferrule* [left=Application]
{\Gamma;\Delta \vdash t : \Pi x : A.\,B \\ \Gamma;\Delta \vdash u : A}
{\Gamma;\Delta \vdash t\,u : B[x/u]}

\inferrule* [left=Abstraction]
{\Gamma;\Delta \vdash A : Type \\ \Gamma;\Delta(x : A) \vdash t : B \\ B \neq Kind}
{\Gamma;\Delta \vdash \lambda x : A.\,t : \Pi x : A.\,B}

\inferrule* [left=Product]
{\Gamma;\Delta \vdash A : Type \\ \Gamma;\Delta(x : A) \vdash B : s}
{\Gamma;\Delta \vdash \Pi x : A.\,B : s}

\inferrule* [left=Conversion]
{\Gamma;\Delta \vdash t : A \\ \Gamma;\Delta \vdash B : s \\ A \equiv_{\beta\Gamma} B}
{\Gamma;\Delta \vdash t : B}
\end{mathpar}
\end{center}
\caption{Typing rules for terms in the $\lambda\Pi$-calculus Modulo (Figure~3 in~\cite{Saillard}).}
\label{fig:lambdapi-figure3}
\end{figure}

The backend encoding in Lambdapi (\texttt{encode.lp}, about 600 lines) defines signatures, constructor-list well-formedness, subtyping checks, and signature-directed elimination on top of the $\lambda \Pi$ calculus modulo rewriting~\cite{Saillard}; it is a shallow encoding of the rules from Section~\ref{sec:rules}. 
The frontend transpiler parses a restricted surface syntax and lowers it to backend terms. In fact, it transpiles only a fragment of the minimal language proposed, but this is enough to demonstrate the simplicity and expressiveness of the system.
The lowering mirrors Rules~1--7 from Section~\ref{sec:rules}. Signature annotations become \texttt{signature\_term}. Constructor introductions become \texttt{signature\_intro}. Signature-only matches become \texttt{signature\_case}. Appendix~\ref{app:transpiler} shows representative input and generated output.

\section{Conclusion}
\label{sec:conclusion}
In this research, we introduced a compact yet powerful subtype system that demonstrates remarkable versatility and compatibility with various type theories. Our system effectively addresses the challenge of reusing and subtyping constructors in indexed datatypes through the novel introduction of first-class datatype signatures. A particularly notable feature of our approach is its ability to eliminate trivial proof obligations even when an arbitrary set of constructors is involved, significantly reducing the proof burden in practical applications. We have illustrated the system's flexibility through several practical examples drawn from real-world programming scenarios. Furthermore, we have provided a concise implementation that maintains simplicity without sacrificing expressiveness, making our approach accessible for integration into existing proof assistants and type systems.

\section{Future Work}
There are several directions in which we would like to extend this work.

\begin{itemize}
\item \emph{Mechanized meta-theory.} We plan to formalize the typing and subtyping rules, together with subject reduction and progress, in a proof assistant such as Coq or Agda, and connect the proofs to our prototype implementation.
\item \emph{Larger case studies.} Our current examples are small; applying constructor subset signatures to larger developments in existing proof assistants (e.g.\ libraries of lists and vectors) would help validate their practical benefits.
\item \emph{Richer signatures.} Finally, we intend to explore simple extensions of signatures with basic set operations on constructor sets (such as unions) while keeping the meta-theory and implementation lightweight.
\end{itemize}

\appendix

\section{Surface-to-Lambdapi Transpiler}
\label{app:transpiler}
The current artifact is organized around a translation pipeline: users write a compact surface fragment, and \texttt{siglp-transpiler} lowers it to backend terms consumed by \texttt{subtyping.encode}. This keeps examples short while preserving explicit generated code.

The supported surface fragment translates LambdaPi \texttt{Inductive} declarations. We compile the \texttt{Def} declarations, signature shorthand (\texttt{\{T :: c\_1 | \dots | c\_n\}}), and signature-only \texttt{match} expressions. For each inductive declaration, the transpiler generates helper symbols for target, index shape, telescope, and constructor list. It rewrites term-level signatures to \texttt{signature\_term}. It rewrites constructor applications to \texttt{signature\_intro}. It rewrites signature matches to \texttt{signature\_case}.

\begin{figure}[!ht]
\begin{lstlisting}[language=ML, basicstyle=\ttfamily\small, keywordstyle=\bfseries\color{blue}, commentstyle=\itshape\color{gray}, stringstyle=\color{red}, numbers=left, numberstyle=\tiny\color{gray}, frame=single]
Static nat : *.
Static 0 : nat.
Static +1 : Nat -> nat.
Nat :: {nat :: |0 |+1}.

Static vec : * -> nat -> *.
Static empty : (s : *) -> (vec s 0).
Static cons : (s : *) -> (n : nat) -> (vec s n) -> (lift s) -> (vec s (+1 n)).

NonEmptyVecNat |v :: nat -> * => {(vec nat (+1 v)) :: |cons}.

head_nonempty |v m :: nat -> (NonEmptyVecNat v) -> Nat =>
[m of Nat
  |(cons k) => (+1 0)
].
\end{lstlisting}
\caption{Mini-language rendering of the non-empty vector example(\emph{The syntax was slightly modified to match the implementation of the minimal language in Haskell})}
\label{fig:transpiler-surface}
\end{figure}

Compiling with \texttt{cabal run siglp-transpiler} (input: \texttt{examples/sample.siglp}, output: \texttt{examples/sample.lp}) yields LambdaPi code such as Figure~\ref{fig:transpiler-output}.

\begin{figure}[!ht]
\begin{lstlisting}[language=Lambdapi, basicstyle={\ttfamily\scriptsize\upshape}, inputencoding=utf8, extendedchars=true]
require open subtyping.encode;

symbol vec__sig_target : Set;
rule vec__sig_target ↪ vecLabel;
symbol vec__sig_indexes : indexes_type;
rule vec__sig_indexes ↪
  index_poly_succ (index_dependent_lift_succ natLabel (index_null));
symbol vec__sig_delta : delta_tel;
rule vec__sig_delta ↪
  mk_delta_tel vec__sig_indexes vec__sig_pattern_any;

symbol head_nonempty : Π (v : nat),
  (signature_term vec__sig_target
    (mk_delta_tel vec__sig_indexes (...))
    (constructor_list_cons consLabel (constructor_list_nil))) → nat;

rule head_nonempty ↪ λ (v : nat) (m : ...),
  (signature_case vec__sig_target (mk_delta_tel vec__sig_indexes (...))
    (constructor_list_cons consLabel (constructor_list_nil))
    natLabel m
    (case_branches_cons consLabel (constructor_list_nil) natLabel
      (λ (k : constructor_type consLabel), (+1 _0))
      (case_branches_nil natLabel)));
\end{lstlisting}
\caption{Generated LambdaPi excerpt after lowering signatures and matches}
\label{fig:transpiler-output}
\end{figure}

The complete generated files are omitted for space. Full examples are available in \texttt{subindex/subtyping/transpiler/examples/}, and the target encoding used by generated terms is \texttt{subindex/subtyping/encode.lp}.

\paragraph*{Backend shallow encoding (\texttt{encode.lp}).}
The backend is organized as a shallow encoding of Rules~1--7 into LambdaPi symbols and rewrite rules. Formation and well-formedness are captured by \texttt{signature\_wfb} and \texttt{signature}. Subtyping uses boolean checkers (\texttt{signature\_subtypeb}, \texttt{signature\_to\_base\_subtypeb}) wrapped by \texttt{is\_true} when a proposition is required. Elimination is represented directly by \texttt{signature\_term}, \texttt{case\_branches}, and \texttt{signature\_case}. Finally, a coercion hook (\texttt{coerce\_rule}) inserts signature widening automatically when constructor-list inclusion is validated.

\begin{figure}[!ht]
\begin{lstlisting}[language=Lambdapi, basicstyle={\ttfamily\scriptsize\upshape}, inputencoding=utf8, extendedchars=true]
symbol signature_wfb : Set → delta_tel → constructor_list → Bool;
inductive signature : Set → delta_tel → constructor_list → TYPE ≔
| signature_nil : Π (T : Set) (g : delta_tel), signature T g constructor_list_nil
| signature_cons :
    Π (T : Set) (g : delta_tel) (c : labelConstructor) (cs : constructor_list),
    π (is_true (against_delta_b T g (constructor c))) →
    signature T g cs →
    signature T g (constructor_list_cons c cs);

symbol signature_subtypeb :
  Set → delta_tel → constructor_list →
  Set → delta_tel → constructor_list → Bool;
  
inductive signature_term : Set → delta_tel → constructor_list → TYPE ≔
| signature_intro :
    Π (T : Set) (g : delta_tel) (phi : constructor_list) (c : labelConstructor),
    π (is_true (constructor_in_signature_b T g phi c)) →
    constructor_type c →
    signature_term T g phi;

symbol signature_case :
  Π (T : Set) (g : delta_tel) (phi : constructor_list) (Q : Set),
  signature_term T g phi →
  case_branches phi Q →
  lift Q;

coerce_rule coerce_signature
  (signature $T $g $phi1)
  (signature $T $g $phi2)
  $x ↪ coerce_signature $T $g $phi1 $phi2 $x;
\end{lstlisting}
\caption{Shallow backend encoding excerpt from \texttt{encode.lp}}
\label{fig:encode-shallow}
\end{figure}

\subsection{IA aid}
Although the transpiler was made with IA aid, the mini language was before the popularization of the tool, and lambdapi encoding was done manually.

\bibliography{lipics-v2021-sample-article}

\end{document}
